\begin{table}[t]
\centering
\small
\begin{threeparttable}
\caption{Structural configurations, mechanism representation, and diagnostic roles of the three investigated model structures.}
\label{tab:structural_configurations}
\begin{tabular}{lccc}
\toprule
Property / Characteristic & Base & TGD & CN \\
\midrule
Shared host parameters & 15 & 15 & 15 \\
Calibrated parameters (shared + added) & 15 + 0 & 15 + 2 & 15 + 2 \\
Temperature used by added structural component & No & Yes & Yes \\
Additional generic temperature-conditioned storage / memory & No & Yes & No \\
Explicit snow accumulation–melt sequence & No & No & Yes \\
Primary diagnostic role & Snow-process omission configuration & Temperature-conditioned generic control & Explicit snow accumulation–melt representation \\
\bottomrule
\end{tabular}
\begin{tablenotes}[flushleft]
\footnotesize
\item \textit{Note}: Base represents the snow-process omission configuration lacking explicit snow dynamics (15 core host parameters). TGD is the parameter-count-matched, temperature-conditioned generic storage control that provides non-specific thermal retention/memory without snow physics (15 host + 2 storage/smoothing parameters: $\tau_{\mathrm{warm}}, \Delta\tau_{\mathrm{cold}}$). CN incorporates an explicit degree-day snow accumulation--melt sequence comprising precipitation phase partitioning, persistent snowpack storage, and temperature-threshold melt release (15 host + 2 snow parameters: snowpack thermal inertia coefficient $C_{\mathrm{TG}}$ and degree-day melt factor $K_f$). Matching the number of added parameters does not imply that TGD and CN encode the same structural information. TGD serves as a generic control rather than an intermediate step in a physical decomposition or an additive performance ladder. CN is the generating structure in the controlled synthetic experiment only; it is not treated as hydrological truth in real catchments. Independent calibration (IC-CMA-ES; catchment-wise optimization) and differentiable parameter learning (dPL-MLP; shared cross-catchment parameter mapping) represent contrasting parameter-estimation constraints evaluated in parallel, not competitive ranking benchmarks.
\end{tablenotes}
\end{threeparttable}
\end{table}
